% !TEX root = ../../tree_notes.tex
% =====================================================================
%  7.2  Insertion and Deletion                 sec:c05-insert-delete
%  Exports: def:c05-insert  thm:c05-delete  prop:c05-preserve
%           lst:c05-inscode  lst:c05-delcode  ex:c05-del-run
%           prob:c05-delete  lst:c05-delete  (practice: LC 701)
% =====================================================================
\section{Insertion and Deletion}\label{sec:c05-insert-delete}

\begin{previouslybox}
\begin{enumerate}[nosep, leftmargin=1.6em]
  \item Where does a search for an \emph{absent} key stop, and what is the last node it looked at?
        (\propref{c05-search}, \coderef{c05-search})
  \item Which subtree contains the inorder \emph{successor} of a node with two children, and how do
        you reach it in one walk? (\thmref{c05-inorder})
  \item Why does the inorder sequence of a BST stay sorted after we edit it -- what must any edit
        preserve? (\defref{c05-prop}, \thmref{c05-inorder})
  \item On $B$ (\figref{c05-running}), which node is the smallest key larger than $3$?
\end{enumerate}
\end{previouslybox}

\begin{topicintro}
A sorted array is easy to \emph{read} and painful to \emph{edit}: inserting costs a $\bigO{n}$ shift.
A BST is the opposite bargain -- editing is just a search with a consequence. \term{Insertion} is a
failed search that hangs the new key where the null link was: one descent, no shifting, the invariant
holds by construction. \term{Deletion} is the chapter's real mechanic, because removing a node must
splice the tree back together \emph{without} breaking $\text{left}<\text{node}<\text{right}$. There
are three cases by child-count, and only the nasty one -- two children -- needs the classic trick:
overwrite the doomed node with its \term{inorder successor} (the smallest key that is still larger),
then delete that successor, which is guaranteed to be a simple node. This section proves the trick
keeps the sorted line intact (\propref{c05-preserve}) and works it twice on $B$.
\textbf{Prerequisites.} The search descent of \secref{c05-property} (\propref{c05-search}) and the
inorder-sorted equivalence (\thmref{c05-inorder}); the \code{TreeNode} pointers of \chref{c02}.
Deletion here is \emph{structural editing}; keeping the height bounded so these edits stay
$\bigO{\log n}$ is Chapter~8 (\chref{c06}).

\textbf{After this section you can}
\begin{itemize}[nosep]
  \item insert a key with one descent that stops exactly at the correct null link; \hfill\lv{2}
  \item delete a node in all three child-cases, including the two-child successor splice; \hfill\lv{3}--\lv{4}
  \item prove that insert and delete leave the inorder sequence sorted (\propref{c05-preserve}); \hfill\lv{3}
  \item decide successor-vs-predecessor replacement and see it does not change correctness. \hfill\lv{3}
\end{itemize}
\end{topicintro}

\begin{recognitionbox}[Pattern card -- edit a BST and keep it a BST]
\reckey{the task adds or removes a key from a tree that must \emph{remain} a valid BST: ``insert'',
``delete'', ``prune'', ``remove one occurrence''.}
\recuse{
  \item insert a value -- \defref{c05-insert}, a search that hangs a leaf (\probref{c05-delete}'s twin
        LC 701 in the practice set);
  \item delete a value -- the three cases and the successor splice (\thmref{c05-delete}, \probref{c05-delete});
  \item ``remove all keys outside $[L,R]$'' (trim) or ``every key in a range'' -- repeated deletion,
        the descent prunes whole subtrees (referenced to \chref{c24}, not re-linked);
  \item any edit that must leave the \emph{sorted line} valid -- check with \propref{c05-preserve}.
}
\recnotuse{
  \item you only \emph{read} (search/count/kth) -- no structural change, \secref{c05-property} or \secref{c05-queries};
  \item the container handles insert/delete for you -- \code{std::set}, \secref{c05-contests};
  \item you must \emph{guarantee} $\bigO{\log n}$ worst case -- that is a \emph{balanced} tree, \chref{c06}.
}
One-line test: \emph{does a pointer need to change so the tree grows or shrinks by one key, while
staying ordered? Then edit it here -- and prove \thmref{c05-inorder} survives.}
\end{recognitionbox}

\subsection{Intuition}

\textbf{Insertion} is the least scary edit. Run the search for the new key $x$; you will bottom out at
a null link (that is what ``absent'' means in \propref{c05-search}). The node you were standing on
when you stepped into that null is $x$'s parent: if $x$ is smaller it is the left child, else the
right. Because $x$ was placed exactly where a comparison sent it, every ancestor's bound is
respected, so the invariant is \emph{born correct} -- no re-balancing, no re-sorting.

\textbf{Deletion} splits by how many children the doomed node $v$ has, because the number of
sub-threads to re-tie is exactly the number of children.

\begin{itemize}[nosep]
  \item \textbf{Leaf (0 children):} detach $v$ from its parent. Done.
  \item \textbf{One child:} bypass $v$ -- connect $v$'s parent straight to $v$'s child. The whole
        subtree of $v$ is already on the correct side of the parent (it was, before $v$ left), so the
        invariant is untouched.
  \item \textbf{Two children:} you cannot detach or bypass -- $v$ holds \emph{two} live subtrees that
        both need a single common parent. So do not remove the \emph{node}; remove a \emph{key}.
        Replace $v$'s key by its inorder successor $s$ (the smallest key larger than $v$), and then
        delete $s$, which lives in $v$'s right subtree and \emph{cannot have a left child} (it was the
        left-most), so it falls to case 1 or 2.
\end{itemize}

\begin{figure}[H]
\centering
\begin{tikzpicture}[diagram,
    del/.style={boxnode, draw=warningred, text=warningred, circle},
    suc/.style={boxnode, draw=successgreen, text=successgreen, circle}]
  % two-child case mini-tree
  \node[del] (v)  at (0,0)   {$v$};
  \node[tnode] (l) at (-2,-1.4) {$L$};
  \node[tnode] (r) at (2,-1.4)  {$R$};
  \node[suc] (s)  at (1,-2.8)  {$s$};
  \node[tnode] (rr) at (3,-2.8) {$R'$};
  \draw (v) -- (l); \draw (v) -- (r);
  \draw (r) -- (s); \draw (r) -- (rr);
  \node[annot, align=left, font=\scriptsize] at (5.2,-1.4)
    {$s=$ left-most of $R$\\ (inorder successor)\\ $s$ has no left child\\ $\Rightarrow$ delete $s$\\ is case 1 or 2};
  \node[annot, align=center, font=\scriptsize] at (-4.4,-1.4)
    {copy $s$'s key into $v$,\\ then delete $s$};
\end{tikzpicture}
\caption{The two-child deletion. Node $v$ (red) has two subtrees $L,R$; we cannot drop or bypass it.
Instead we take $s$ -- the inorder successor, the left-most node of $R$ (green) -- copy $s$'s key into
$v$, and delete the \emph{original} $s$. Because $s$ was left-most it has no left child, so deleting
$s$ is the trivial leaf/one-child case. Every key stays on its correct side: $s>v$ and every $L$ key
$<s$, so the sorted line is preserved (\propref{c05-preserve}).}\label{fig:c05-del2}
\end{figure}

Choosing the \emph{successor} (largest of the small / smallest of the large) rather than the
\emph{predecessor} is a convention, not a law: symmetrically you could copy the inorder predecessor
(the largest key smaller than $v$) out of $v$'s left subtree and delete that. Either way, the node
you actually unlink has at most one child.

\subsection{Definitions and Notation}

\begin{definition}{Insertion as a failed search}{c05-insert}
To insert key $x$ into BST $B$: run the search descent for $x$ until the current pointer is null;
let $p$ be the last non-null node visited. Attach a new leaf carrying $x$ as $p$'s left child if
$x<p$'s key, else as $p$'s right child. If $B$ was empty, $x$ becomes the root. The result satisfies
\defref{c05-prop} and adds no other structure.
\end{definition}

\noindent\textbf{In plain words.} Insert \emph{is} search, stopping one step earlier: instead of
reporting ``absent,'' you plant the key at the null spot you hit. It costs the same $\bigO{h}$ and
never disturbs any existing pointer except one null that becomes a leaf.

\begin{theorem}{Deletion: three cases and the successor splice}{c05-delete}
To delete a key $k$ present at node $v$ in a BST: (i) if $v$ has no children, unlink it; (ii) if $v$
has exactly one child, replace $v$ by that child; (iii) if $v$ has two children, let $s$ be $v$'s
inorder successor (left-most node of $v$'s right subtree), copy $s$'s key into $v$, and delete $s$
from the right subtree (which is then a case (i)/(ii) node). The output is a BST with key $k$ removed
and every other key retained.
\end{theorem}
\begin{proofbox}
Cases (i) and (ii) only remove or re-attach whole subtrees whose keys already formed a legal interval
under $v$'s parent, so \defref{c05-prop} is preserved at the parent. In (iii): $s$ is the smallest
key $>v$'s key, so every key formerly $<v$'s key stays $<s$ (they are $<v<s$), and every other key
formerly $>v$'s key stays $>s$ \emph{except} $s$ itself, which we then remove. Copying $s$ into $v$
therefore leaves $v$'s left subtree all $<s$ and its right subtree (minus $s$) all $>s$. As $s$ is
left-most in the right subtree it has no left child, so its own removal is case (i) or (ii). The two
keys ``$v$ and $s$'' collapse to one surviving node holding $s$'s key; the multiset of keys loses
exactly $k$. \end{proofbox}
\begin{insightbox}
Deletion never ``moves'' a subtree wholesale -- it moves a \emph{single value} so that the node we
actually unlink is a simple one. That is the whole trick, and why LC 450 rates medium: the idea is a
sentence, the pointer bookkeeping is the interview.
\end{insightbox}

\begin{proposition}{Insert and delete preserve the sorted line}{c05-preserve}
If $B$ is a BST with strictly increasing inorder, then after either \defref{c05-insert} of a new key
or \thmref{c05-delete} of an existing key, the resulting tree is a BST: its inorder is again
strictly increasing. Insert appends one element to the sorted line; delete removes one.
\end{proposition}
\begin{proofbox}
By \thmref{c05-inorder} it suffices to show the inorder stays increasing. Insert places $x$ as a leaf
under $p$ with $x<\text{key}(p)$ on the left (so $x$ lands between $p$'s predecessor and $p$ in
inorder) or $x>\text{key}(p)$ on the right (so it lands just after $p$); every ancestor's bound held
by the very comparisons of the descent, so the sequence is still sorted. Delete (case iii) swaps
$v$'s key to $s$ where $s$ was already $v$'s inorder successor -- adjacent in the sorted line -- so
erasing $s$ and overwriting $v$ is exactly ``remove $k$ from a sorted list''; cases (i)/(ii) remove a
node with no ordering consequence for the rest. Hence sortedness survives. \end{proofbox}

\subsection{The algorithm: plant a leaf, or splice a value}

\begin{lstlisting}[caption={Insert a key into a BST -- a search that stops and plants (LeetCode 701)}, label={lst:c05-inscode}]
TreeNode* insertIntoBST(TreeNode* root, int val) {
    if (root == nullptr) return new TreeNode(val);      // empty tree: the leaf IS the root
    TreeNode* cur = root;
    for (;;) {
        if (val < cur->val) {
            if (cur->left) cur = cur->left;             // keep descending left
            else { cur->left  = new TreeNode(val); break; }  // null hit: plant here
        } else {
            if (cur->right) cur = cur->right;           // keep descending right
            else { cur->right = new TreeNode(val); break; }
        }
    }
    return root;   // same root as before -- insertion never changes the root pointer
}
\end{lstlisting}

\begin{lstlisting}[caption={Delete a key from a BST -- the three cases with successor splice (LeetCode 450)}, label={lst:c05-delcode}]
// Smallest-key node of a (non-empty) subtree, and its parent, in one left-walk.
TreeNode* leftmost(TreeNode* node) { while (node->left) node = node->left; return node; }

TreeNode* deleteNode(TreeNode* root, int key) {
    if (root == nullptr) return nullptr;              // key absent: nothing to do
    if (key < root->val) {
        root->left = deleteNode(root->left, key);     // descend left
    } else if (key > root->val) {
        root->right = deleteNode(root->right, key);   // descend right
    } else {                                          // found the node to delete
        if (root->left == nullptr) {                  // 0 or 1 child (only right)
            TreeNode* up = root->right; delete root; return up;
        }                                             // bypass to the surviving child
        if (root->right == nullptr) {                 // 1 child (only left)
            TreeNode* up = root->left; delete root; return up;
        }
        // two children: copy successor key down, then delete the successor
        TreeNode* succ = leftmost(root->right);
        root->val = succ->val;
        root->right = deleteNode(root->right, succ->val);  // succ has no left child -> easy case
    }
    return root;
}
\end{lstlisting}
\begin{complexitybox}
Both edits cost $\bigO{h}$: insertion is one descent; deletion is one descent to $v$ plus (case iii)
one descent into the right subtree to find (and remove) the successor. Space is $\bigO{h}$ for the
recursive deletion's call stack ($\bigO{1}$ for the iterative insert). As always the honest cost is
$h$, not $\log n$ -- a chain makes every edit $\bigO{n}$, the rot \secref{c05-contests} warns about
and \chref{c06} cures.
\end{complexitybox}

\subsection{Worked Example}

\begin{example}{Deleting $3$ and then $6$ from the running BST $B$}{c05-del-run}
\emph{Delete $3$ (two children $1,6$):} $3$'s inorder successor is the left-most node of its right
subtree, i.e.\ the min under $6$ -- that is $4$. Copy $4$ into the node that held $3$, then delete the
original leaf $4$ (it was $6$'s left child, a case (i) leaf). New tree: node now reading $4$ has left
child $1$ and right child $6$ (whose children are now just $7$); inorder
$1,4,6,7,8,10,13,14$ -- sorted, $3$ gone. \checkmark\\
\emph{Delete $6$ (from the original $B$, two children $4,7$):} $6$'s successor is the left-most of its
right subtree $=\{7\}$, so $s=7$. Copy $7$ into $6$'s node, then delete the leaf $7$ (case (i)). Now
$6$'s old node holds $7$ with a single left child $4$; inorder $1,3,4,7,8,10,13,14$ -- sorted, $6$
gone. \checkmark In both, the node we physically unlinked was a leaf, exactly as \thmref{c05-delete}
guarantees.
\end{example}
\begin{tracebox}[Why the successor is always safe to splice]
Delete $3$: successor $s=4$. Is $4$'s right-subtree role legal at node $3$? After copying $4$ into
$3$'s node, we need every left key $<4$ (left child is $1$: ok) and every right key $>4$ (right
subtree rooted $6$, all $\ge 6$: ok). And $4$ must still fit node $3$'s \emph{ancestors}: it is under
$8$'s left, needs $<8$ -- $4<8$ \checkmark. Then deleting the leaf $4$ is case (i). The successor,
being left-most, could never have a left child, so we never recurse into a fresh two-child case --
one splice and we are done. This is \propref{c05-preserve} working as promised.
\end{tracebox}

\subsection{Variations and Twists}

\begin{itemize}[nosep]
  \item \textbf{Predecessor variant}: replace with the right-most node of the left subtree instead of
        the successor; mirror all the inequalities. Same complexity, different (but valid) shape. \hfill\lv{3}
  \item \textbf{Recursive insert}: \code{if (!root) return new TreeNode(val);} recurse into one child
        and return \code{root}. Elegant, $\bigO{h}$ stack; the loop version is stack-safe. \hfill\lv{2}--\lv{3}
  \item \textbf{Duplicate keys}: adopt ``duplicates go right'' and make insert/delete count them, or
        store a multiplicity in each node so insert of an existing key bumps a counter and delete
        decrements before unlinking. \hfill\lv{3}
  \item \textbf{Bulk delete (trim)}: to remove all keys $<L$ you descend once; when a node is $<L$ its
        entire left subtree goes and you re-attach its right -- $\bigO{h + \text{removed}}$. Trim BST
        is the named exercise (\chref{c24}), not re-linked here. \hfill\lv{3}--\lv{4}
  \item \textbf{Parent tracking}: an iterative delete needs the parent of $v$ to re-point it; the
        recursion sidesteps this by returning the replacement to the caller -- one of several reasons
        LC 450 is cleaner recursively. \hfill\lv{3}
\end{itemize}

\subsection{Common Pitfalls}\label{sec:c05-delete-pitfalls}

\begin{warningbox}[Trap 1 -- in the two-child case, splicing the successor \emph{node} out of its old
place by hand]
Beginners unlink the successor node from the right subtree with fiddly ``if it has a right child,
graft it'' code and forget the successor might have a \emph{right} child (it can -- it only lacks a
\emph{left} child). The clean way is recursion: \code{root->right = deleteNode(root->right, succ->val)}
lets the same three-case routine delete the successor, now a case (i)/(ii) node. Reimplementing the
splice inline is where the bugs live.
\end{warningbox}
\begin{lstlisting}[caption={\textcolor{warningred}{WRONG}: hand-splice successor, ignoring its possible right child}, language=C++, basicstyle=\ttfamily\small]
root->val = succ->val;
root->right->left = succ->right;   // breaks when succ is root->right itself (no left link!)
\end{lstlisting}
\begin{lstlisting}[caption={\textcolor{successgreen}{RIGHT}: reuse the routine to delete the successor key}, language=C++, basicstyle=\ttfamily\small]
root->val = succ->val;
root->right = deleteNode(root->right, succ->val);   // succ is left-most -> easy case, safely handled
\end{lstlisting}

\begin{warningbox}[Trap 2 -- forgetting the key is absent]
If \code{key} is not in the tree the routine must return the tree unchanged. The base case
\code{if (!root) return nullptr;} covers it: descending for an absent key bottoms out at a null, and
rebuilding the pointers on the way up is a no-op. Testing only present keys hides this; LC 450
explicitly allows absent keys.
\end{warningbox}

\begin{warningbox}[Trap 3 -- leaking or dangling the removed node]
Deleting node $A$ while a local \code{up} still aliases it, or \code{delete}-ing then reading $A$,
is undefined. Capture the replacement pointer \emph{before} freeing, as in \coderef{c05-delcode}. In
LeetCode's managed environment you may skip \code{delete}, but the pointer order still matters for
correctness on the two-child case.
\end{warningbox}

\subsection{Problem Lab}

The section's anchor is LC 450 -- delete a node from a BST, keeping it a BST. It is the single
problem that forces every idea here at once: the descent to find the node, the child-count case split,
and the successor splice with clean recursive bookkeeping. It is medium precisely because the \emph{idea}
is one line and the \emph{pointers} are the test.

% ---------------------------------------------------------------------------
\solvedproblem{c05-delete}{Delete Node in a BST}{LeetCode 450}
  {https://leetcode.com/problems/delete-node-in-a-bst/}{LC Medium L3}
  {BST delete three cases, inorder successor splice, pointer reattachment, recursive return}

\begin{problemstatement}[Delete Node in a BST]
Given a \code{root} node of a BST and a \code{key}, \textbf{delete} the node whose value equals
\code{key} and return the root of the \emph{modified} BST. The deletion must \emph{maintain the BST
property}. The key may not be present (return the original tree). $0\le n\le 10^{4}$ nodes, values in
$[-10^{4},10^{4}]$ and distinct.
\end{problemstatement}

\begin{attemptgate}[Attempt it yourself first -- \lv{3}, give yourself 25 minutes]
This is the chapter's pointer-mechanic test. The two-child case is where attempts stall.
\peek{How many children does a node have to let you just unlink it? What can you do when a node has
two live subtrees that both need one parent? Think \emph{key}, not \emph{node}.}
\nudge{Descend to the key. If it has $0$ or $1$ child, reconnect its parent to its surviving child.
If it has two, find its inorder successor (left-most of the right subtree), \emph{copy its value} into
the doomed node, and then \emph{delete the successor} -- which now has at most one child, so it drops
into the easy case.}
\fullhint{Recursive: \code{key<root->val} recurse left and assign; \code{key>root->val} recurse
right; equal $\Rightarrow$ handle three cases. Two children: \code{succ} = left-most of
\code{root->right}; \code{root->val = succ->val}; \code{root->right = deleteNode(root->right,
succ->val)}. Return \code{root}. \coderef{c05-delcode} is the full build.}
\end{attemptgate}

\solstep{Understanding the Problem}
We must remove one key and hand back a still-valid BST. The statement's promise that the result is a
BST is the whole difficulty: unlinking a node is easy, but re-tying pointers so no ordering is broken
-- especially when the node has two subtrees -- is the task. Absence of the key is allowed, so our
routine must be a no-op then. The expected cost is $\bigO{h}$: one descent plus (two-child case) one
descent to find the successor.

\solstep{Key Observations}
\begin{observation}
The number of \emph{children} decides the mechanics: $0$ children $\Rightarrow$ detach; $1$ child
$\Rightarrow$ bypass; $2$ children $\Rightarrow$ you cannot move two subtrees under one parent, so move
a single \emph{value} instead (\thmref{c05-delete}).
\end{observation}
\begin{observation}
The inorder successor is always in the right subtree and always \emph{left-most}, so it has no left
child -- deleting it is automatically a case $0$/$1$ problem. This makes the recursion terminate
without a second two-child tangle.
\end{observation}
\begin{observation}
Copying the successor's key down and then deleting the successor is exactly ``remove $k$ from the
sorted inorder line'': \propref{c05-preserve} guarantees the rest stays sorted.
\end{observation}

\solstep{Think It Through With Me}
\begin{thinkbox}
\deadend{Physically detach a two-child node by grafting its left subtree under some node of the right
subtree (or vice-versa).}
{You can do it -- find the right-most of the left subtree and hang the right subtree there -- but now
you have written a second, more delicate splice, and its height can grow. The successor-copy trick
turns the hard node into an easy one in \emph{one} descent, and the easy node is deleted by the very
routine you already wrote. Simpler and provably order-preserving.}
\deadend{Search the whole tree for the key and rebuild a fresh BST without it.}
{Correct but $\bigO{n}$ and ignores the structure. The descent finds the node in $\bigO{h}$; there is
no reason to touch subtrees the comparisons already proved irrelevant (\propref{c05-search}).}
\aha{Descend to the key; branch on child-count. For two children, overwrite the node with its inorder
successor's value and recursively delete the successor (which has $\le 1$ child). One routine handles
all three cases because the successor's deletion recurses into the easy branch.}
\end{thinkbox}

\solstep{Solution}
Recurse on \code{root}. Base: null returns null. If \code{key < root->val}, re-point
\code{root->left} at the recursive delete of \code{key}; if \code{key > root->val}, the symmetric right call;
otherwise \code{root} is the target: if it lacks a left child, return its right child; if it lacks a
right child, return its left child; else find the left-most node of the right subtree, copy its value
into \code{root->val}, and recursively delete that successor key from \code{root->right}. Return
\code{root}. Each returned pointer is stitched into the caller's child slot, so bypasses and detach
propagate correctly, and the top-level \code{root} is the (unchanged or replaced) tree root.

\solstep{Proof of Correctness}
\begin{proofbox}
Induction on the descent depth; the recursion's outer calls only reassign a child to the delete
result of that same child, so it suffices to prove correctness when \code{root} is the target.
\emph{Case left null:} \code{root} has $\le 1$ (right) child; every key in that child's subtree
exceeds the deleted key and lies within \code{root}'s parent interval, so returning it re-attaches a
legal subtree (\propref{c05-preserve}). Symmetric for right null. \emph{Two children:} let $s$ be the
left-most key of the right subtree -- the smallest key $>$ the deleted key. Copying $s$ into
\code{root} keeps \code{root}'s left subtree all $<s$ (they are $<\text{deleted}<s$) and its right
subtree (minus $s$) all $>s$. By induction the recursive \code{deleteNode(root->right, s)} returns the
right subtree with $s$ removed and still a BST; assigning it to \code{root->right} yields a valid BST
whose key multiset is the original minus the deleted key. The successor has no left child, so its own
deletion is a null-child case and the recursion does not add a new two-child level. Absent key: the
comparisons route to a null base case and rebuild pointers unchanged. Hence in all cases the returned
tree is a valid BST with key removed (or unchanged if absent). \end{proofbox}

\solstep{Implementation}
\begin{lstlisting}[caption={Delete Node in a BST -- submittable solution with successor splice (LeetCode 450)}, label={lst:c05-delete}]
/**
 * Definition for a binary tree node.
 * struct TreeNode { int val; TreeNode *left, *right; };
 */
class Solution {
    TreeNode* leftmost(TreeNode* n) { while (n->left) n = n->left; return n; }
public:
    TreeNode* deleteNode(TreeNode* root, int key) {
        if (!root) return nullptr;                       // absent -> unchanged
        if (key < root->val) {
            root->left  = deleteNode(root->left,  key);  // one descent, one child
        } else if (key > root->val) {
            root->right = deleteNode(root->right, key);
        } else if (!root->left) {                        // target: 0 or 1 child
            TreeNode* up = root->right; delete root; return up;
        } else if (!root->right) {                       // target: 1 child (left)
            TreeNode* up = root->left;  delete root; return up;
        } else {                                         // target: two children
            TreeNode* succ = leftmost(root->right);      // inorder successor, no left child
            root->val    = succ->val;
            root->right  = deleteNode(root->right, succ->val);
        }
        return root;
    }
};
\end{lstlisting}

\solstep{Trace}
\begin{tracebox}[Running BST $B$: delete $3$ (the two-child path), then delete an absent key]
\emph{delete(3):} descend $8\to3$ (match). $3$ has left $1$ and right $6$ $\Rightarrow$ two-child case.
\code{leftmost(3->right)} walks $6\to4$ (then $4$'s left is null) $\Rightarrow$ \code{succ=4}. Set
\code{3->val=4}. Recurse \code{delete(3->right=6,\ 4)}: $4<6$ go left to $4$; match, $4$ has no
children so its parent $6$ gets \code{6->left = nullptr} and the leaf is freed. Tree now: $8$ left
child $4$(was $3$) with children $1$ and $6$; $6$ has only right child $7$. \checkmark\\
\emph{delete(5):} $5<8\to3(\text{now }4)$? compare with $4$: $5>4$ go right to $6$; $5<6$ go left to
null $\Rightarrow$ base case returns null, $6$'s left stays null, pointers rebuilt unchanged. Tree
identical. \checkmark (absent key is a no-op, Trap 2).
\end{tracebox}

\solstep{Complexity}
\begin{complexitybox}
Time $\bigO{h}$: one descent to the target, plus one left-descent into its right subtree for the
successor (and the successor's deletion is the $\bigO{1}$-child branch). Space $\bigO{h}$ recursion
stack. Worst case $h=n$ on a chain; balanced $h=\bigO{\log n}$ -- the height caveat of
\secref{c05-property} again, and the reason \chref{c06} balances. Verdict: \ac.
\end{complexitybox}

\solstep{Pitfalls and Variants}
\begin{warningbox}
Three classic failures: (i) hand-splicing the successor's parent link (Trap 1) -- \code{succ} may sit
directly at \code{root->right}, so ``\code{parent->left = succ->right}'' writes a link that does not
exist; recursion avoids the special case entirely. (ii) Deleting a present node but forgetting to
reassign \code{root->left/right}, so the child pointer still names freed memory. (iii) Treating the
two-child case as ``just move the whole subtree,'' which changes the shape and can rebalance height
unboundedly. \emph{Variants:} the predecessor mirror (replace with right-most of the left subtree);
and trim-by-range, where the same descent deletes many keys with $\bigO{h+\text{removed}}$.
\end{warningbox}

\begin{insightbox}[Key takeaways]
\begin{enumerate}[nosep, leftmargin=1.6em]
  \item \textbf{Move a value, not a subtree.} The two-child case is scary only until you realise the
        successor is a \emph{simple} node whose key can fill the hole (\thmref{c05-delete}).
  \item \textbf{Recursion is the splice.} Returning the replacement into the caller's child slot
        makes detach, bypass and successor-removal one code path -- the cleanest fix for Trap 1.
  \item \textbf{Ordering is inherited, not enforced.} Because the descent followed the comparisons,
        every edit lands where \defref{c05-prop} already holds -- that is \propref{c05-preserve}.
\end{enumerate}
\end{insightbox}

\begin{hardeningbox}
\begin{itemize}[nosep, leftmargin=1.4em]
  \item \textbf{Iterative delete} needs the parent of both the target and the successor; the returned
        child must be grafted onto \code{parent}. Doable, but three extra pointers -- one reason the
        recursive form dominates in interviews.
  \item \textbf{Randomised delete}: some balanced BSTs merge the two subtrees by a random choice
        between successor and predecessor to keep shapes varied -- \chref{c06} territory (treaps).
  \item \textbf{Delete-and-reinsert idiom}: for ``remove keys in range,'' repeatedly \code{deleteNode}
        each offending key -- correct but $\bigO{k\cdot h}$; a single pruned recursion does it in
        $\bigO{h+\text{removed}}$ (the trim exercise, \chref{c24}).
\end{itemize}
\end{hardeningbox}

\subsection{Practice Problems}

\begin{stuckbox}
\begin{itemize}[nosep, leftmargin=1.4em]
  \item Did your insert \emph{stop at the null link} the search found, not re-search from the root?
        (\defref{c05-insert}.)
  \item In the two-child delete, did you copy the successor's \emph{value} and then delete the
        successor, rather than unlinking a subtree? (\thmref{c05-delete}, Trap 1.)
  \item Does an \emph{absent} key leave the tree untouched? (Trap 2, base case of \coderef{c05-delcode}.)
  \item After your edit, is the inorder still strictly increasing? (\propref{c05-preserve}.)
\end{itemize}
\end{stuckbox}

\begin{practiceset}
  \practice{Insert into a Binary Search Tree}{LeetCode 701}
    {https://leetcode.com/problems/insert-into-a-binary-search-tree/}{LC Medium L3}
    {Given a BST (possibly empty) and a value \emph{not present} in it, insert it and return the root,
    keeping the BST valid. This is \defref{c05-insert} verbatim -- run the descent to the null link and
    plant the leaf; the ``medium'' rating only exists because you must keep a \code{parent} pointer
    (or recurse and assign the returned node into the child slot) so the new leaf actually attaches.
    Any correct insertion position is accepted -- you do not have to reproduce a unique shape, which
    is exactly why the one-descent \coderef{c05-inscode} earns \ac.}
\end{practiceset}

\begin{recallbox}[10-second refresher -- 7.2 Insertion and Deletion]
Insert $=$ a search that \emph{plants} at the null link it bottoms out on (\defref{c05-insert},
\coderef{c05-inscode}): one descent, born correct. Delete splits by child-count (\thmref{c05-delete}):
leaf $\Rightarrow$ detach; one child $\Rightarrow$ bypass; two children $\Rightarrow$ copy the inorder
\emph{successor}'s value (left-most of the right subtree, so it has no left child) and recursively
delete that successor (\coderef{c05-delcode}). Both edits keep the inorder sorted -- \propref{c05-preserve}.
Cost is $\bigO{h}$ each; the successor splice, not the idea, is the interview difficulty.
\selftest{
  \item Give the exact null link where you attach key $5$ to $B$ (\figref{c05-running}) and say $5$'s parent.
  \item Why can the inorder successor of a two-child node never have a left child?
  \item Delete $8$ (the root) from $B$: name the successor, where it is copied, and which node is physically removed.
  \item State how \propref{c05-preserve} follows from \thmref{c05-inorder} for a two-child deletion.
  \item Write the three base cases of \coderef{c05-delcode} from memory and say which one the successor lands in.
}
\end{recallbox}
